\documentclass[11pt]{article}
\usepackage[letterpaper,margin=0.85in]{geometry}
\usepackage{amsmath,amssymb,amsthm}
\usepackage{xcolor}
\usepackage[colorlinks=true,linkcolor=blue,urlcolor=blue]{hyperref}
\setlength{\parindent}{0pt}
\setlength{\parskip}{6pt}
\setlength{\emergencystretch}{2em}

\newtheorem{prop}{Proposition}
\newtheorem{cor}{Corollary}

\title{G1a: Two-Parameter Scale Existence\\
\large Radius and length cannot be stabilized by the naive positive-energy model}
\author{}
\date{2026-08-07 07:32 CST (UTC+8)}

\begin{document}
\maketitle

\section{Question and verdict}

The one-parameter restriction establishes that conserved axial tube flux and
adiabatic plasma content prevent collapse under an \emph{isotropic} rescaling. G1a
asks the stronger question: does the same energy possess a finite stationary point
when tube radius and total network length can change independently?

\begin{center}
\fbox{\parbox{0.9\textwidth}{
\textbf{Verdict.} No. In the stated thin-tube model, every point that minimizes the
length direction has a strictly downhill radius direction. The isotropic minimum is
a constrained minimum, not a two-parameter equilibrium. Axial flux and adiabatic
plasma content prevent point collapse but do not fix the network's aspect or tube
radius. A second invariant or geometric/field-energy term is required.
}}
\end{center}

This is a model-level falsification, not a no-go theorem for all graph reactors.

\section{Model and scaling assumptions}

Let $r>0$ scale the common branch radius and $L>0$ scale every branch length of a
fixed theta-graph shape. Absorb reference geometry and sums over edges into positive
coefficients. The energy proposed for the first G1a test is
\begin{equation}
E(r,L)=
AL+DrL+Fr^2L+B\frac{L}{r^2}
+C(r^2L)^{-2/3},
\label{eq:energy}
\end{equation}
where:
\begin{itemize}
\item $AL$ is an effective line-tension cost;
\item $DrL$ is tube interface or surface energy;
\item $Fr^2L$ is a core or condensate volume cost;
\item $BL/r^2$ is axial magnetic energy with tube flux frozen; and
\item $C(r^2L)^{-2/3}$ is the internal energy of a fixed-content adiabatic plasma
with $\gamma=5/3$ and volume proportional to $r^2L$.
\end{itemize}

All coefficients are nonnegative; for the strict result below, axial flux gives
$B>0$. The model assumes slender tubes, a common radius scaling, fixed graph shape,
and no separate poloidal/twist flux, external pressure, bending energy, mutual
inductance, or hard geometric constraint.

\section{Exact reduction}

Define
\[
g(r)=A+Dr+Fr^2+\frac{B}{r^2},
\qquad
E(r,L)=Lg(r)+Cr^{-4/3}L^{-2/3}.
\]
For each fixed $r$, $E\to\infty$ as $L\to0$ and $L\to\infty$, and the unique
length minimizer follows from $\partial_LE=0$:
\begin{equation}
L_*(r)=
\left[
\frac{2C}{3g(r)}r^{-4/3}
\right]^{3/5}.
\label{eq:Lstar}
\end{equation}
At this point the plasma energy is $3/2$ times $L_*g$, so the exactly reduced energy
is
\begin{equation}
E_*(r)=\frac52\left(\frac{2C}{3}\right)^{3/5}
\left[\frac{g(r)}{r^2}\right]^{2/5}.
\label{eq:Ered}
\end{equation}

\begin{prop}[Two-parameter G1a no-stationary-point result]
If $A,D,B>0$ and $F\ge0$, the energy in Eq.~\eqref{eq:energy} has no stationary
point at finite $(r,L)\in(0,\infty)^2$.
\end{prop}

\begin{proof}
The only radius-dependent factor in Eq.~\eqref{eq:Ered} is
\[
H(r)=\frac{g(r)}{r^2}
=F+\frac{D}{r}+\frac{A}{r^2}+\frac{B}{r^4}.
\]
Its derivative is
\[
H'(r)=-\frac{D}{r^2}-\frac{2A}{r^3}-\frac{4B}{r^5}<0
\]
for every finite $r>0$. Therefore $E_*(r)$ is strictly decreasing and cannot
possess a finite-radius stationary point. Since every full stationary point would
also be stationary after exact minimization in $L$, none exists.
\end{proof}

The same obstruction appears directly. At $\partial_LE=0$, write the plasma term
$P=C(r^2L)^{-2/3}$; then $P=\tfrac32Lg(r)$. Substitution into
$r\partial_rE$ gives
\[
r\partial_rE
=-L\left(2A+Dr+\frac{4B}{r^2}\right)<0.
\]
The core-volume coefficient $F$ cancels: a positive energy depending only on total
volume cannot choose the aspect of that volume.

\section{Why the isotropic pass was insufficient}

Under the restricted deformation $r=\lambda r_0$, $L=\lambda L_0$, the energy has
positive powers and negative powers of $\lambda$ and can possess the unique minimum
of the restricted energy. But the full plane contains an anisotropic direction in which $r$
increases while $L$ decreases. If $F>0$, the minimizing trajectory approaches
constant volume ($L_*\sim r^{-2}$) while line, surface, and axial-flux costs fall;
the infimum is reached only at $r\to\infty$, $L\to0$. This runaway eventually
violates the slender-tube approximation, which is precisely the signal that the
minimal model does not close its own scale problem.

\section{Scientific consequence}

G1a has failed for the naive energy but succeeded as a gate: it found the missing
physics. A viable model needs at least one independent shape-sensitive contribution.
Candidates include:
\begin{itemize}
\item conserved poloidal or twist flux, whose magnetic energy can scale as $1/L$;
\item helicity and the full cycle-flux/mutual-inductance matrix of the theta graph;
\item bending or curvature energy and junction-core energy;
\item an external-pressure or vessel constraint, stated rather than hidden; or
\item a relaxed-MHD spectral constraint in a genus-2 domain.
\end{itemize}

The next G1a calculation should add the least speculative of these---conserved axial
and poloidal fluxes---and determine whether the resulting finite critical point has a
positive $2\times2$ Hessian without relying on the boundary of the thin-tube regime.

\section{Verification and scope}

Every identity above is independently checked in
\path{../code/g1a_two_parameter_checks.py}. The script also displays the isotropic
restriction and verifies the direct stationarity contradiction symbolically.

This result does not address junction topology, branch kink/sausage stability,
reconnection, transport, fusion gain, or engineering. It disproves only finite-scale
equilibrium for Eq.~\eqref{eq:energy} under the declared independent scalings.



\end{document}
